\begin{center}\begin{tikzpicture}[font=\small,>=stealth,line join=round]
\draw[gray!40,rounded corners=3pt] (0,0) rectangle (14.2,8.3);\node at (7.1,7.700000000000001) {镜子在1与3：半格出发，两次反射后到四点五};
\begin{scope}[shift={(1.6901639344262296,3.814705882352941)},x=0.9836065573770493cm,y=1.352941176470588cm]\draw[->,gray] (-0.6,0)--(5.5,0) node[right]{$x$};\draw[->,gray] (0,-1.6)--(0,1.8) node[above]{$y$};\draw[blue,thick,solid,domain=-0.2:5,samples=180] plot(\x,{cos(deg(pi*(\x-1)/2))});\fill[orange!85!black] (0.5,0.7071067811865476) circle(2.2pt);\fill[orange!85!black] (1.5,0.7071067811865476) circle(2.2pt);\fill[orange!85!black] (4.5,0.7071067811865476) circle(2.2pt);\foreach \a in {1,3}{\draw[violet,dashed] (\a,-1.3)--(\a,1.3);\node[above] at (\a,1.3){$x=\a$};}\draw[orange!85!black,dashed] (.5,.70710678)--(4.5,.70710678);\end{scope}\node[anchor=west,align=left,text width=5.0cm] at (8.4,4.0) {$f=\cos\frac{\pi(x-1)}2$\\[9pt]轴 $x=1,3$\\[9pt]$0.5\to1.5\to4.5$\\[9pt]三点同高 $\frac{\sqrt2}2$\\[9pt]横移4，两倍轴间距};
\node[text width=12.6cm,align=center] at (7.1,.65) {图中反射点来自曲线；任意起点同样满足两次反射后水平移动4。};
\end{tikzpicture}\end{center}
